% !TEX program = xelatex
\documentclass[11pt]{ctexart}
%=====================================================================
% theme-violet.tex —— 讲义精讲 · 紫花主题
% 用法：\documentclass[11pt]{ctexart}
%       \input{assets/theme-violet}
%       \begin{document} ... \end{document}
% 只用 ctexart 自带能力 + xcolor/enumitem/fancyhdr/pifont/amsmath，
% 不依赖 tikz/tcolorbox，保证 MiKTeX 可直接编译。
%=====================================================================

\usepackage[a4paper,left=2.0cm,right=2.0cm,top=2.1cm,bottom=2.3cm,%
            headsep=12pt,footskip=22pt]{geometry}
\usepackage{amsmath,amssymb}
\usepackage{xcolor}
\usepackage{array}
\usepackage{enumitem}
\usepackage{pifont}
\usepackage{fancyhdr}
\usepackage[hidelinks]{hyperref}

%---------------------------------------------------------------- 配色
\definecolor{violet}{RGB}{107,78,143}      % 主色：紫
\definecolor{violetdeep}{RGB}{63,44,92}    % 深紫：标题/强调
\definecolor{petal}{RGB}{183,159,208}      % 花瓣：装饰线、花饰
\definecolor{bandviolet}{RGB}{247,243,251} % 浅紫底：档案/提示框
\definecolor{bandcream}{RGB}{252,249,242}  % 米白底：易错/注意框
\definecolor{rulepurple}{RGB}{217,207,230} % 分隔线
\definecolor{inkgray}{RGB}{92,96,110}      % 次要文字

%---------------------------------------------------------------- 字体
% 有 Noto CJK 用 Noto，没有就退回 ctex 默认（Windows 下为宋体/黑体）
\IfFontExistsTF{Noto Serif CJK SC}{\setCJKmainfont{Noto Serif CJK SC}}{}
\IfFontExistsTF{Noto Sans CJK SC}{\setCJKsansfont{Noto Sans CJK SC}}{}

%---------------------------------------------------------------- 版式
\setlist{nosep,leftmargin=1.9em}
\linespread{1.28}
\setlength{\parskip}{3.5pt}
\setlength{\parindent}{0pt}
\setlength{\tabcolsep}{5pt}
\renewcommand{\arraystretch}{1.28}

%---------------------------------------------------------------- 花饰
% \ding{95} = 五瓣小花，\ding{96} = 星形花；都来自 pifont，必定存在
\newcommand{\orn}{\textcolor{petal}{\ding{95}}}          % 行内小标
\newcommand{\ornstar}{\textcolor{petal}{\ding{96}}}      % 星形花
\newcommand{\flowerrule}{%                                % 居中花饰分隔线
  \par\vspace{7pt}%
  \noindent\makebox[\linewidth]{\color{petal}%
    \rule{0.20\linewidth}{0.6pt}\hspace{8pt}{\small\ding{95}}\hspace{4pt}%
    \raisebox{1.1pt}{\tiny\ding{95}}\hspace{4pt}{\small\ding{95}}\hspace{8pt}%
    \rule{0.20\linewidth}{0.6pt}}%
  \par\vspace{7pt}}

%---------------------------------------------------------------- 标题与节头
\newcommand{\doctitle}[3]{%   #1 主标题  #2 副标题  #3 元信息（可用 \\ 分行）
  \begin{center}
    {\color{petal}\small \orn\enspace #2\enspace \ornstar}\\[7pt]
    {\LARGE\bfseries\color{violetdeep} #1}\\[10pt]
    {\color{petal}\rule{0.55\linewidth}{1pt}}\\[4pt]
    {\small\color{inkgray} #3}
  \end{center}\vspace{2pt}}

\newcommand{\hdr}[2]{%   #1 章节号（如 §0）  #2 标题
  \par\vspace{12pt}%
  \noindent\colorbox{violetdeep}{\color{white}\bfseries\small\;#1\;}\hspace{7pt}%
  {\color{violetdeep}\large\bfseries #2}\par\nobreak
  \vspace{3pt}{\color{petal}\hrule height 1.2pt}\vspace{6pt}}

\newcommand{\subhdr}[1]{%
  \par\vspace{6pt}%
  \noindent{\color{violet}\bfseries #1}\par\vspace{1pt}}

%---------------------------------------------------------------- 提示框
\newcommand{\callout}[1]{% 浅紫底：讲义档案、读法说明
  \par\medskip\noindent
  \fcolorbox{rulepurple}{bandviolet}{\parbox{\dimexpr\linewidth-2\fboxsep-2\fboxrule}{\vspace{2pt}#1\vspace{2pt}}}%
  \par\medskip}

\newcommand{\warnbox}[1]{% 米白底：易错、边界条件
  \par\medskip\noindent
  \fcolorbox{rulepurple}{bandcream}{\parbox{\dimexpr\linewidth-2\fboxsep-2\fboxrule}{\vspace{2pt}#1\vspace{2pt}}}%
  \par\medskip}

\newcommand{\gapitem}[2]{% #1 缺口类型（紫） #2 内容
  \par\vspace{2pt}\noindent
  \hangindent=1.4em\hangafter=1
  {\color{petal}\small\ding{96}}\enspace{\bfseries\color{violetdeep}#1}\enspace #2\par}

% 前提旁注：只放在「需要理解的那一步」后面，给陈述 + 本题代入，不重证。
\newcommand{\aside}[1]{%
  \par\vspace{1pt}\noindent
  \hangindent=1.6em\hangafter=1
  {\small\color{inkgray}%
    \hspace*{0.6em}{\color{petal}$\triangleright$}\hspace{2pt}#1}%
  \par\vspace{2pt}}

%---------------------------------------------------------------- 表格
% 本机 MiKTeX 的 colortbl 与 array 版本不匹配（\rowcolor 会报错），
% 故表格统一用「上粗紫线 + 表头 + 中细线」的三线式，不依赖 colortbl。
\newcommand{\tabtop}{\par\vspace{3pt}{\color{violet}\hrule height 1.2pt}\vspace{1pt}}
\newcommand{\tabmid}{\par\vspace{1pt}{\color{petal}\hrule height 0.5pt}\vspace{1pt}}
\newcommand{\tabbot}{\par\vspace{3pt}{\color{rulepurple}\hrule height 0.5pt}\vspace{8pt}}
\newcommand{\thead}[1]{{\bfseries\color{violetdeep}#1}}

%---------------------------------------------------------------- 页眉页脚
\pagestyle{fancy}
\fancyhf{}
\renewcommand{\headrulewidth}{0.4pt}
\renewcommand{\headrule}{\hbox to\headwidth{\color{rulepurple}\leaders\hrule height \headrulewidth\hfill}}
\fancyhead[L]{\small\color{inkgray}\songti \doctag}
\fancyhead[R]{\small\color{inkgray}\songti \coursetag}
\fancyfoot[C]{\small\color{inkgray}\orn\enspace{\color{violet}\thepage}\enspace\orn}

% 每篇讲义在导言区设这两个值
\newcommand{\doctag}{}
\newcommand{\coursetag}{}
\newcommand{\settags}[2]{\renewcommand{\doctag}{#1}\renewcommand{\coursetag}{#2}}

%---------------------------------------------------------------- 数学记号
\newcommand{\vect}[1]{\boldsymbol{#1}}
\newcommand{\dd}{\,\mathrm{d}}
\newcommand{\ddx}{\frac{\mathrm{d}}{\mathrm{d}x}}
\addtolength{\jot}{2pt}



\settags{降阶与 Abel 公式 · 精讲}{常微分方程 · 轻松通关手账}

\begin{document}

\doctitle{已知一解求通解：降阶法与 Abel 公式}{常微分方程 · 核心方法精讲}{依据课件 7(1).pdf（L7）p30--p36 与 8(1).pdf（L8）p50--p53，逐页读文字层\quad｜\quad 整理 2026-10-21}

\callout{\subhdr{\orn\enspace 讲义档案与问题坐标}
\textbf{讲什么}：只给一个非零特解 $y_p$，怎么把二阶线性齐次方程的通解写出来。\\
\textbf{记号}：$y_p$ 表示已知的那个解（课件写成 $y_p$，有些书写成 $y_1$）；$y_2$ 表示求出来的第二个解；$W$ 为朗斯基行列式。\\
\textbf{依据}：L7 p30 的式 (17)(18) 给的是\textbf{替换}，L8 p50 的 Theorem 给的是\textbf{朗斯基公式}。两讲之间隔了一节课，课件\textbf{没有}把它们合成一条公式——本讲义做的就是这个拼接，并在 \S10 逐条列出补了哪些缺口。}

\hdr{§0}{先回答“老师讲了吗”}
讲了，而且是两处、两块：

\tabtop
\begin{tabular}{@{}p{0.30\linewidth}p{0.64\linewidth}@{}}
\thead{位置} & \thead{内容}\\
\end{tabular}
\tabmid
\begin{tabular}{@{}p{0.30\linewidth}p{0.64\linewidth}@{}}
L7 p30，式 (17)(18) & 一般结论：$n$ 阶线性齐次方程，用 $y=y_p\int z\dd x$ 可把阶数降一\\
L7 p31--p33，例 1 & $x^3y''+(x-1)y'-y=0$，$y_p=x-1$\\
\textbf{L7 p33--p36，例 2} & $y''+xy'-(x+1)y=0$，$y_p=e^x$\\
L8 p50，Theorem & Abel 公式 $W=W(x_0)\exp\left(-\int_{x_0}^{x}a_1\dd t\right)$\\
L8 p51--p53，例 & $2xy''-(x+4)y'+\left(1+\frac4x\right)y=0$，由 $y_1=x$ 求 $y_2$\\
L7 p46，作业 III & $xy''+2y'-xy=0$，$y_p=\dfrac{e^x}{x}$（另一题，别与例 2 混）\\
\end{tabular}
\tabbot

\subhdr{三点要说清楚：}
\begin{enumerate}
\item \textbf{课件例 2 的答案就是积分形式}：$y=e^x\left[C_5+C_4\int\exp\left(-\frac{x^2}{2}-2x\right)\dd x\right]$（p36）。它和公式写法完全一致，不是“没算完”——这个积分本来就不是初等函数（\S6 给证明）。
\item \textbf{课件只走“代回方程”这一条路}。Abel 那条“一步写出通解”的合并公式，课件没有单列；但两块原料（降阶替换 + Abel 定理）都在，本讲义把它们接起来，并证明两条路线给出的\textbf{是同一个积分}（\S4）。
\item \textbf{作业不是例 2}。作业 III 是 $xy''+2y'-xy=0$、$y_p=e^x/x$。两题长得像（都是变系数、都给一个指数型特解），直接套例 2 的中间结果必错。
\end{enumerate}

\hdr{§1}{降阶为什么成立（课件 p30 没讲的那一步）}
课件 p30 直接写出式 (18) 的替换，一句话“the order is reduced by unity”。为什么？本节的展开就是那个“为什么”。

设 $y=y_p(x)\,v(x)$，其中 $v$ 是待求函数。求导：
\[
y'=y_p'v+y_pv',\qquad
y''=y_p''v+2y_p'v'+y_pv'' .
\]
代入算子 $L[y]=y''+a_1(x)y'+a_2(x)y$，按 $v'',v',v$ 三堆归并：
\[
L[y]=y_p\,v''+\left(2y_p'+a_1y_p\right)v'+v\cdot\underbrace{\left(y_p''+a_1y_p'+a_2y_p\right)}_{=\,L[y_p]} .
\]
因为 $y_p$ 是解，$L[y_p]\equiv0$，最后一项整体消失：
\begin{equation}\label{eq:core}
L[y]=y_p\left[\,v''+\left(\frac{2y_p'}{y_p}+a_1\right)v'\,\right].
\end{equation}
在 $y_p\neq0$ 的区间上，$L[y]=0$ 等价于方括号为零。方程里\textbf{已经没有 $v$ 本身，只剩 $v'$}——这就是“降阶”的全部机制。

\warnbox{\subhdr{这不是“试出来的”}\;$v$ 项消失不是巧合：$L$ 是线性算子，把 $y_pv$ 展开后 $v$ 的系数必然正是 $L[y_p]$，而它恒为零。任何线性齐次方程、任何阶数都成立，课件式 (18) 只是这个事实的二阶写法。}

\hdr{§2}{路线一：把方括号里的方程解出来}
记 $z=v'$，式 \eqref{eq:core} 变成一阶线性齐次方程
\[
z'+\left(\frac{2y_p'}{y_p}+a_1\right)z=0 .
\]
分离变量（$z\not\equiv0$ 时）：
\[
\frac{z'}z=-\left(\frac{2y_p'}{y_p}+a_1\right)
\ \Longrightarrow\
\ln|z|=-2\ln|y_p|-\int a_1\dd x+C ,
\]
指数上正好除出一个 $y_p^2$：
\begin{equation}\label{eq:z}
z=C_1\,\frac{\exp\left(-\int a_1\dd x\right)}{y_p^{2}} .
\end{equation}
再积一次分，并不定积分里的常数写进 $C_2$：
\[
v=C_2+C_1\int\frac{\exp\left(-\int a_1\dd x\right)}{y_p^{2}}\dd x,
\qquad
y=y_p\,v .
\]
于是通解
\begin{equation}\label{eq:general}
y=C_1\,y_p+C_2\,y_p\int\frac{\exp\left(-\int a_1\dd x\right)}{y_p^{2}}\dd x .
\end{equation}
课件在例 1、例 2 里没有写这条通用公式，而是每题现场分离变量（因为 $\frac{2y_p'}{y_p}+a_1$ 恰好简单）。写成 \eqref{eq:general} 就是“公式版”。

\hdr{§3}{路线二：Abel 公式一步到位}
先把方程化成\textbf{标准形}（首项系数为 $1$）：
\[
y''+a_1(x)y'+a_2(x)y=0 .
\]
两个解的朗斯基行列式定义为 $W=y_py_2'-y_p'y_2$。Abel 定理（L8 p50 的式 (43)）说：对区间 $I$ 上的任意 $x_0\in I$，
\begin{equation}\label{eq:abel}
W(x)=W(x_0)\exp\left(-\int_{x_0}^{x}a_1(t)\dd t\right).
\end{equation}
关键的一步是\textbf{把 $W$ 除以 $y_p^2$}，左边正好凑成商的导数：
\[
\left(\frac{y_2}{y_p}\right)'=\frac{y_py_2'-y_p'y_2}{y_p^{2}}=\frac{W}{y_p^{2}} .
\]
把 \eqref{eq:abel} 代进去，再积分：
\[
y_2=y_p\int\frac{W(x_0)\exp\left(-\int a_1\dd x\right)}{y_p^{2}}\dd x+C\,y_p .
\]
$C\,y_p$ 是齐次解，并进 $C_1$ 即可；$W(x_0)$ 只是个常数，取 $1$（理由见 \S4）。得到
\begin{equation}\label{eq:abel-sol}
y_2=y_p\int\frac{\exp\left(-\int a_1\dd x\right)}{y_p^{2}}\dd x .
\end{equation}
\textbf{这个方法的特点}：只看 $a_1(x)$ 和 $y_p(x)$，$a_2(x)$ 完全不出现——因为 Abel 公式里只有 $a_1$。这也是它在考场上快的原因。

\hdr{§4}{两条路线是同一个积分（课件缺的拼接）}
对比 \eqref{eq:z} 与 \eqref{eq:abel-sol}：
\begin{itemize}
\item 路线一解出的一阶方程的解：$z=C\,\dfrac{\exp\left(-\int a_1\right)}{y_p^{2}}$；
\item 路线二直接用的被积函数：$\dfrac{W}{y_p^{2}}=W(x_0)\cdot\dfrac{\exp\left(-\int a_1\right)}{y_p^{2}}$。
\end{itemize}
两者只差常数 $W(x_0)$；这个常数积分一次后变成 $C\,y_p$，被通解的 $C_1$ 吸收。所以：

\callout{\textbf{Abel 法不是另一个方法}，它只是把路线一里“解那个一阶方程”的一步，换成了一个现成公式。两条路线的积分核完全相同——这就是为什么两讲的内容可以合成一条公式：
\[
\boxed{\;y_2=y_p\int\frac{W}{y_p^{2}}\dd x,\qquad W=\exp\left(-\int a_1\dd x\right)\;}
\]}

考场上推荐路线二（快、不容易在代数上翻车），检查时用路线一（把 $y_2$ 代回原方程验算）。

\hdr{§5}{例 2 全程：$y''+xy'-(x+1)y=0$，$y_p=e^x$}
先验证特解：$y_p'=y_p''=e^x$，代入得
\[
e^x+x\,e^x-(x+1)e^x=e^x\left(1+x-x-1\right)=0 .\ \checkmark
\]
设 $y=e^x v$，则
\[
y'=e^x(v'+v),\qquad y''=e^x\left(v''+2v'+v\right).
\]
代入并约去 $e^x\neq0$：
\[
\left(v''+2v'+v\right)+x\left(v'+v\right)-(x+1)v=0 .
\]
\textbf{逐项抵消}（课件 p34--p35 只写了一句“the terms containing $v$ cancel”，这里给出明细）：
\[
\underbrace{v}_{\text{来自 }y''}+\underbrace{xv}_{\text{来自 }xy'}-\underbrace{xv}_{\text{来自 }-(x+1)y}-\underbrace{v}_{\text{来自 }-(x+1)y}=0 ,
\]
四个 $v$ 项两两归零，剩下
\[
v''+(x+2)v'=0 .
\]
令 $z=v'$：
\[
z'+(x+2)z=0\ \Longrightarrow\ z=C_4\exp\left(-\frac{x^{2}}2-2x\right).
\]

\warnbox{\subhdr{课件 p35--p36 的一个逻辑缝}\;p35 先假设 $z\neq0$，得 $C_3\neq0$；p36 才补一句“it is easy to see that $z\equiv0$ is a solution”，把常数放开成 $C_4\in\mathbb{R}$。为什么可以补？\\
$z\equiv0\iff v'\equiv0\iff v\equiv\text{常数}\iff y=C\,y_p$——它不是什么新解，\textbf{正是通解里 $C_1y_p$ 那一支}。所以把 $z$ 的常数从 $\neq0$ 放开到 $\mathbb{R}$，既没有多出解，也没有丢解。}

再积一次：
\[
v=C_5+C_4\int\exp\left(-\frac{x^{2}}2-2x\right)\dd x,
\qquad
y=e^x\left[C_5+C_4\int\exp\left(-\frac{x^{2}}2-2x\right)\dd x\right].
\]
这与课件 p36 的答案逐字一致。\textbf{路线二（三行）}：$a_1=x$，故
\[
W=\exp\left(-\int x\dd x\right)=e^{-x^{2}/2},
\qquad
y_2=e^x\int\frac{e^{-x^{2}/2}}{e^{2x}}\dd x=e^x\int e^{-x^{2}/2-2x}\dd x .
\]
与路线一同一结果。

\hdr{§6}{那个积分为什么积不出来}
配平方：
\[
-\frac{x^{2}}2-2x=-\frac{(x+2)^{2}}2+2 ,
\]
所以
\[
\int\exp\left(-\frac{x^{2}}2-2x\right)\dd x
=e^{2}\sqrt{\frac{\pi}{2}}\ \operatorname{erf}\!\left(\frac{x+2}{\sqrt2}\right)+C .
\]
$\operatorname{erf}$ 是误差函数，\textbf{不是初等函数}。这解释了为什么课件、以及任何标准答案都把它留在积分号里——不是偷懒，是这个积分积不出来。考试写到本节上面那一行就是完整答案。

\hdr{§7}{另外两个例子：一眼看清两条路线的分工}
\subhdr{例 1（L7 p31--p33）：}$x^3y''+(x-1)y'-y=0$，$y_p=x-1$。验证：$y_p'=1$、$y_p''=0$，代入 $(x-1)-(x-1)=0$ ✓。\\
\textbf{路线二}：化标准形 $\Rightarrow a_1=\dfrac{x-1}{x^{3}}$，$\displaystyle\int a_1\dd x=-\frac1x+\frac1{2x^{2}}$，故
\[
W=\exp\left(\frac1x-\frac{1}{2x^{2}}\right),
\qquad
z=\frac{W}{y_p^{2}}=\frac{\exp\left(\frac1x-\frac{1}{2x^{2}}\right)}{(x-1)^{2}} .
\]
一步就得到课件 p33 那串分离变量的结果，中间那行 $x^3(x-1)z'+\left[2x^3+(x-1)^2\right]z=0$ 都不用写。通解（$x\neq0,1$）：
\[
y=(x-1)\left[C_1+C_2\int\frac{\exp\left(\frac1x-\frac{1}{2x^{2}}\right)}{(x-1)^{2}}\dd x\right].
\]

\subhdr{例 3（L8 p51--p53）：}$2xy''-(x+4)y'+\left(1+\frac4x\right)y=0$，$y_1=x$。化标准形：
\[
a_1=-\frac{x+4}{2x},\qquad
\int a_1\dd x=-\frac x2-2\ln|x|\ \Longrightarrow\ W\propto x^{2}e^{x/2}.
\]
于是
\[
\frac{y_2}{y_1}=\int\frac{x^{2}e^{x/2}}{x^{2}}\dd x=\int e^{x/2}\dd x=2e^{x/2},
\qquad
y_2=x\,e^{x/2}.
\]
通解 $y=C_1x+C_2x\,e^{x/2}$，与课件 p53 一致。这个例子值得记：$W$ 里的 $x^2$ 与 $y_1^2$ 正好约掉——\textbf{很多题就是设计成“$W$ 与 $y_p^2$ 大量约分”}，看到复杂系数先别慌，先算 $W/y_p^2$ 再化简。

\hdr{§8}{公式卡：四步用法}
\tabtop
\begin{tabular}{@{}p{0.08\linewidth}p{0.46\linewidth}p{0.40\linewidth}@{}}
\thead{\#} & \thead{做什么} & \thead{注意}\\
\end{tabular}
\tabmid
\begin{tabular}{@{}p{0.08\linewidth}p{0.46\linewidth}p{0.40\linewidth}@{}}
1 & 化成标准形 $y''+a_1y'+a_2y=0$ & 首项系数必须是 1，这步错则全错\\
2 & 确认 $y_p$ 真的是解（代一下，30 秒） & 题目给了也要验，防抄错题\\
3 & 算 $W=\exp\left(-\int a_1\dd x\right)$ & 只差常数，取 1；能化简先化简\\
4 & $y_2=y_p\int W/y_p^2\dd x$，写 $y=C_1y_p+C_2y_2$ & 积不出来就保留积分形式\\
\end{tabular}
\tabbot

\textbf{收尾检查}：直接算 $W(y_p,y_2)$，看它是否等于第 3 步取的 $W$（差常数也算对）。这一条能抓住绝大多数符号错误。

\hdr{§9}{六条易错点}
\begin{enumerate}
\item \textbf{没化标准形}。例 1 的 $a_1$ 是 $\frac{x-1}{x^3}$，不是 $x-1$；例 3 的是 $-\frac{x+4}{2x}$，不是 $-(x+4)$。
\item \textbf{$W$ 只差常数倍}。例 3 真实的 $W(x,xe^{x/2})=\frac{x^2}{2}e^{x/2}$，而 $\exp\left(-\int a_1\right)=x^2e^{x/2}$，差 $\frac12$——不影响通解，别以为算出 $\frac12$ 就错了。
\item \textbf{$y_p$ 的零点要挖掉}。例 1 的 $y_p=x-1$ 在 $x=1$ 处为零；例 3 还要挖掉 $x=0$（系数奇点）。通解只在 $y_p\neq0$ 且系数连续的区间上成立。
\item \textbf{非初等积分不是错}。例 2 的 $\int e^{-x^2/2-2x}\dd x$ 就是积不出来的（\S6）。
\item \textbf{只降一阶}。$n$ 阶方程用 $y=y_p\int z\dd x$ 得到的是 $n-1$ 阶，不是一阶（课件原话 reduced by unity）。二阶才恰好降成一阶。
\item \textbf{别和参数变易法混}。两处都有 $W$、都有 $\exp\left(-\int a_1\right)$，但位置相反：求\textbf{齐次}第二解时 $W$ 在\textbf{分子}（$y_2=y_p\int W/y_p^2$）；求\textbf{非齐次}特解时 $W$ 在\textbf{分母}（$C_1'=-y_2f/W$）。记成“求通解 $W$ 上，求特解 $W$ 下”。
\end{enumerate}

\flowerrule

\hdr{§10}{本讲义补了哪些缺口}
这一节是给“从头到尾打通逻辑”用的：课件哪里跳了步、本讲义补了什么，逐条对应。

\gapitem{机制缺口}{课件 p30 只给式 (17)(18) 的替换，未说明为什么降阶。\S1 给出 $L[y_pv]$ 的完整展开，指出 $v$ 项系数恰为 $L[y_p]\equiv0$。}

\gapitem{代数跳步}{课件 p34 末写到 $v''+2v'+v+xv'+xv-xv-v=0$，p35 一句“the terms containing $v$ cancel”直接给出 $v''+(x+2)v'=0$。\S5 把四个 $v$ 项逐项配对写出来。}

\gapitem{平凡解}{课件 p35 先设 $z\neq0$，p36 才补“$z\equiv0$ is a solution”，未解释为何能放开常数。\S5 的提示框说明 $z\equiv0\Leftrightarrow y=Cy_p$，正是通解的第一支。}

\gapitem{拼接缺口}{L7 的降阶替换与 L8 的 Abel 定理分属两讲，课件从未合并。\S3--\S4 把两者接成一条公式，并证明积分核相同。}

\gapitem{适用区间}{课件全程默认 $y_p\neq0$，未点明这是定理成立的前提。\S9 第 3 条把三个例子的奇点逐一标出。}

\gapitem{答案形态}{课件把 $\int e^{-x^2/2-2x}\dd x$ 留在积分号里，未说明原因，容易被误读成“没算完”。\S6 用配平方 + $\operatorname{erf}$ 证明它非初等。}

\hdr{§11}{本讲三句话}
\begin{enumerate}
\item 已知一个解 $y_p$，设 $y=y_pv$，含 $v$ 的项必然消失（因为 $L[y_p]=0$），方程降成只含 $v'$ 的一阶方程。
\item 那条一阶方程的解恒等于 $W/y_p^2$，而 $W=\exp\left(-\int a_1\dd x\right)$，所以 $y_2=y_p\int W/y_p^2\dd x$——$a_2$ 从头到尾不出现。
\item 先化标准形再读 $a_1$；$W$ 差常数不要紧；积分积不出来就保留积分形式。这三条能解释本讲 90\% 的“我算得对不对”。
\end{enumerate}

\hdr{§12}{与课件的页码对照}
\noindent\begin{minipage}{\linewidth}
\tabtop
\begin{tabular}{@{}p{0.14\linewidth}p{0.30\linewidth}p{0.48\linewidth}@{}}
\thead{本讲义} & \thead{课件位置} & \thead{内容}\\
\end{tabular}
\tabmid
\begin{tabular}{@{}p{0.14\linewidth}p{0.30\linewidth}p{0.48\linewidth}@{}}
\S1 & L7 p30，式 (17)(18) & 降阶替换（本讲义补了推导）\\
\S2 & L7 p30--p36 各例现场做法 & 路线一：解出那个一阶方程\\
\S3 & L8 p50，式 (43) & Abel 定理 / 朗斯基公式\\
\S4 & 讲义自补 & 两条路线的等价性（课件无）\\
\S5 & L7 p33--p36，例 2 & $y''+xy'-(x+1)y=0$，$y_p=e^x$\\
\S6 & 讲义自补 & 该积分非初等的证明（课件无）\\
\S7 & L7 p31--p33；L8 p51--p53 & 例 1、例 3\\
\S9 & 讲义自补 & 六条易错（含 $y_p\neq0$ 的适用范围）\\
\end{tabular}
\tabbot
\end{minipage}

\vspace{2pt}
{
{\small\color{inkgray}
\textbf{复核记录}：本文公式沿用 2026-10-20 的 SymPy 逐条验证结果（22 项全通过），本次仅重排版式与补写文字，未改动任何结论。\par}
}

\end{document}

